\chapter{绪\hspace{6pt}论}

\section{研究背景与意义}

固井质量评价是油气井建井和长期生产过程中的关键环节。水泥环承担封隔地层、保护套管和维持井筒完整性的作用。当水泥环与套管或地层之间存在胶结不良、微环空或窜槽时，可能造成层间流体连通、产量下降以及井筒完整性风险。因此，利用测井资料对水泥胶结状态进行识别和评价具有明确的工程意义。

声波测井和超声成像测井是固井质量评价中的重要资料来源。声波测井具有覆盖范围广、成本相对较低、适合全井段普查等优势，但其全波列信号受到多种模式波、套管波、地层波和井眼环境的共同影响，传统的低维幅度或能量指标难以充分表达复杂窜槽结构。CAST 等超声成像测井能够提供较高分辨率的环向声阻抗图像，在窜槽位置和方位结构方面具有更直观的信息。若能利用 CAST 图像构造监督信号，并从 XSI 声波全波列中学习窜槽相关结构，将有助于探索声波资料在固井质量评价中的更充分利用。

然而，声波测井与超声成像测井并非天然对齐。二者的测量物理、采样分辨率、深度采样方式和方位参考均存在差异。尤其在本文所关注的垂直井场景中，XSI 与 CAST 的相对方位参考不稳定，旧实验未能支持可靠的点对点方位监督。因此，本文不将 CAST 环向像素直接作为 XSI 方位通道的逐点标签，而是围绕方位失配问题构造弱监督标签，研究声波 CWT 时频特征与 CAST 派生窜槽结构特征之间的可学习关系。

\placeholderfigure{固井质量、水泥环、窜槽通道、套管和地层之间关系的手绘示意。}{固井质量评价与窜槽问题示意图。该图用于说明水泥环窜槽可能导致层间连通和井筒完整性风险，不包含实验性能结论。}{fig:cement-channeling}

\section{固井质量评价与窜槽识别研究现状}

传统固井质量评价通常依赖声波幅度、变密度测井、声波全波列特征以及超声成像资料等信息。声波幅度类指标解释简单，但容易受到套管、地层和仪器条件影响；全波列资料包含更丰富的传播和散射信息，但人工解释和低维特征提取难以完整利用其时频结构。超声成像资料能够提供更直观的环向图像，对识别水泥环局部缺陷具有优势，但其作业成本、测量条件和覆盖范围也限制了其作为唯一评价手段的使用。

近年来，基于机器学习和深度学习的方法被用于测井解释和图像识别任务。对于固井质量评价问题，深度模型可以从高维波形或时频图中提取非线性特征。但模型性能依赖可靠标签，若标签构造没有处理好多源测井之间的空间、深度或方位不一致，就可能得到表面上较高但实质上受泄漏或错配影响的指标。因此，本文将标签工程和样本划分作为方法设计的核心部分，而不是仅仅比较模型主干。

\section{声波测井与超声成像测井方法概述}

XSI 阵列声波测井记录多接收器、多方位通道的全波列响应。该类数据包含不同传播模式在时间和频率上的叠加，适合通过连续小波变换进行时频展开。CAST 超声成像测井则以环向扫描方式得到声阻抗图像，可用于描述水泥胶结状态的方位分布。本文将 CAST 声阻抗阈值化后构造 severity map，并进一步生成一维 percentage label 或方位向 FFT magnitude label。

\placeholderfigure{左侧为 XSI 多通道声波波形或 CWT 时频图，右侧为 CAST 环向 Zc 图像，中间标注二者数据形态和分辨率差异。}{XSI 与 CAST 测井响应差异示意图。该图用于说明声波时频响应与 CAST 环向成像之间的数据形态差异。}{fig:xsi-cast-difference}

\section{深度学习在测井解释中的应用现状}

深度学习模型能够从高维输入中自动提取层级特征，已被广泛用于图像识别、序列分析和地球物理数据解释等任务。对于声波全波列，直接使用原始时间序列可能难以显式刻画瞬态频率变化；连续小波变换将一维信号映射到二维时频域，使卷积神经网络能够处理局部时间--频率结构。EfficientNetV2B0 作为高效卷积网络主干，可用于从 CWT 图像中提取特征并执行回归任务。

需要强调的是，本文不以模型结构复杂度作为主要创新点。本文重点在于：在 XSI 与 CAST 方位失配条件下，如何构造更稳健的 CAST 派生弱监督标签，并用 single-well depth-heldout 实验检验其可学习性。

\section{现有研究存在的问题}

结合旧项目实验与最终证据冻结，本文面对的主要问题包括：

\begin{enumerate}
  \item XSI 与 CAST 的方位参考不可靠，直接点对点方位监督不适合作为主线。
  \item 一维 percentage label 直观但会丢弃方位结构，且标签稀疏时 zero baseline 在 MAE 上很强。
  \item FFT magnitude label 能降低对绝对方位的依赖，但丢弃相位后不再保留完整方位位置信息。
  \item 连续深度样本之间具有相似性，random split 可能带来相邻深度泄漏风险。
  \item 当前最终证据仅覆盖 \texttt{array\_03} 单井 depth-heldout，不支持多井泛化或工业部署结论。
\end{enumerate}

\section{本文研究内容与技术路线}

本文围绕方位失配条件下的声波测井窜槽结构特征反演开展研究。主要内容包括：构建 XSI CWT 时频输入；基于 CAST Zc 构造 severity map、一维 percentage label 和方位向 FFT magnitude label；建立 CWT-EfficientNetV2B0 回归模型；区分 random split 探索实验与 single-well depth-heldout 验证实验；分析 EXP-008 主线结果、EXP-007 fallback 结果及 EXP-006 exploratory baseline。

\placeholderfigure{从 XSI 波形、CWT、CAST Zc、severity map、FFT magnitude label、EfficientNet 回归、depth-heldout 验证到误差分析的流程图。}{本文技术路线图。该图展示本文从数据构建、标签构造到模型训练和结果审计的整体流程。}{fig:technical-route}

\section{本文组织结构}

第 1 章介绍研究背景、意义和本文研究内容。第 2 章说明 XSI 与 CAST 数据基础，并分析方位失配和样本划分问题。第 3 章提出旋转不变标签构建和 CWT-EfficientNetV2B0 回归方法。第 4 章给出 random split 探索实验、EXP-008 主实验、EXP-007 fallback 实验和 EXP-006 baseline 的结果分析。第 5 章讨论标签设计意义、zero baseline MAE 问题、误差来源和单井泛化限制。第 6 章总结全文并给出后续工作方向。
